\section{Introdução}
\label{sec:introducao}

A qualidade interna de \emph{software} descreve propriedades estruturais do código-fonte, como modularidade, coesão e acoplamento, que podem ser investigadas por análise estática \cite{iso25010,pressman:21}. Em sistemas orientados a objetos, a suíte de Chidamber e Kemerer oferece indicadores para caracterizar essas dimensões \cite{chidamber:94}. Esses indicadores não constituem, isoladamente, uma nota de maturidade: sua interpretação depende do tamanho, do domínio e da organização do projeto.

O ecossistema Spring permite investigar a evolução de projetos Java ao longo de diferentes períodos tecnológicos. Projetos iniciados antes da consolidação do Spring Boot podem ter atravessado mudanças de configuração, dependências, práticas de desenvolvimento e ferramentas de assistência. O Spring Boot 1.0 estável foi lançado em 1º de abril de 2014 \cite{springboot:2014}; por isso, o período de 2009 a 2013 é tratado como contexto de Spring tradicional, com identificação explícita de eventuais versões preliminares de Boot. A presença atual de Boot em um repositório criado em 2010 não demonstra que essa tecnologia já integrava seu código em 2010.

A mesma distinção é necessária para a inteligência artificial. Um projeto antigo pode apresentar hoje vestígios de ferramentas de IA introduzidos recentemente. Datas de criação dos repositórios não datam o uso dessas ferramentas, e fronteiras de calendário não separam perfeitamente produção humana e assistida. O Copilot já estava em prévia técnica em junho de 2021 \cite{copilot:2021}. Evidências de arquivos, identidades e mensagens de commit permitem observar parte dessa adoção, mas não todo uso nem a proporção do código produzida por IA \cite{robbes:26}.

A pergunta central deste trabalho é: \textbf{como a qualidade estrutural de projetos do ecossistema Spring evoluiu ao longo do tempo, e quais mudanças acompanharam a adoção do Spring Boot e o aparecimento de evidências de assistência por IA?} O foco está no acompanhamento dos mesmos projetos, especialmente aqueles criados entre 2009 e 2013. As demais coortes oferecem contexto sobre projetos surgidos em épocas posteriores. Esse desenho distingue evolução dentro do projeto de diferenças de composição entre conjuntos de projetos.

O objetivo geral é descrever longitudinalmente essa evolução por versões anuais recuperadas do histórico Git. Os objetivos específicos são: (i) caracterizar trajetórias de acoplamento, complexidade, tamanho, herança e coesão; (ii) descrever mudanças nos mesmos projetos em torno da primeira presença observável de Spring Boot; (iii) descrever mudanças em torno do aparecimento de evidências de IA; e (iv) comparar essas trajetórias com o contexto de projetos sem a mesma transição observada, considerando tamanho, idade e configuração do framework. O paradigma GQM liga essas perguntas às sete medidas do instrumento CK.

A relevância do estudo está em produzir evidências históricas rastreáveis, com código, datas, versões e cobertura associados a cada medição. A mineração de repositórios possibilita esse acompanhamento, mas exige cuidado com seleção, disponibilidade do histórico e interpretação dos metadados \cite{hassan:08,kalliamvakou:16}. Mudanças de framework, linguagem, domínio, equipe e tamanho são explicações concorrentes para variações das métricas. Nem a injeção de dependências elimina a influência do framework sobre a estrutura, nem um contraste entre períodos permite atribuir diferenças exclusivamente ao desenvolvedor ou à IA.

Os resultados incluem tabelas de cobertura, trajetórias anuais, medianas por época e diferenças entre versões do mesmo projeto, acompanhadas de um pacote de reprodução. A investigação admite estabilidade, aumentos e reduções das métricas, sem presumir melhoria. A Seção~\ref{sec:fund} apresenta a fundamentação; a Seção~\ref{sec:trabalhos}, os trabalhos relacionados; a Seção~\ref{sec:metodologia}, o protocolo histórico; e a Seção~\ref{sec:resultados}, a execução histórica e seus limites, distinguindo-os dos pilotos transversais anteriores.
